% !TeX spellcheck = en_us
\documentclass[a4paper,11pt,fleqn]{article}
\usepackage[left=1.50cm, right=1.50cm, top=3cm, bottom=1cm, 
headheight= 2.5cm, headsep= 4pt, footskip=25pt]{geometry}
\usepackage{fancyhdr}
\usepackage{lmodern}
\usepackage[table]{xcolor}
\usepackage{amssymb,amstext,amsmath}
\usepackage{eso-pic}
\usepackage{booktabs}
\usepackage{bm}
\usepackage{setspace}
\usepackage{graphicx}
\usepackage[colorlinks]{hyperref}
\setlength{\parindent}{0mm}
\setlength{\parskip}{2mm}
\usepackage{enumitem}
\usepackage{titlesec}
\usepackage{gensymb}
\usepackage{pgfgantt}
\usepackage{lscape}
\usepackage[absolute]{textpos}
\usepackage{comment}
\usepackage{cite}
\usepackage{tikz}
\usepackage{pgfplots}
\usetikzlibrary{shapes.geometric}
\usepackage{array}
\newcolumntype{P}[1]{>{\centering\arraybackslash}p{#1}}
\usepackage{longtable}
\usepackage[font={small}]{caption}
\usepackage[most]{tcolorbox}
\usepackage{tikz}
\usepackage{stmaryrd}
%\usepackage[none]{hyphenat}
\usepackage{ctable}
\usepackage{ragged2e}
\usepackage[official]{eurosym}
\newcommand*\circled[1]{\tikz[baseline=(char.base)]{
		\node[shape=circle,draw,inner sep=1pt] (char) {#1};}}
\usepackage[symbol]{footmisc}
\renewcommand{\thefootnote}{\fnsymbol{footnote}}
\renewcommand{\thempfootnote}{\fnsymbol{mpfootnote}}
\renewcommand*\footnoterule{}

\usepackage{xcolor}
\definecolor{docColor}{cmyk}{1.00	, 1.00  , 0   , 0.40}  % 0.90\% of black
\color{docColor}

\newtcolorbox{myhlbox}[2][]{colback=white, colframe=docColor, width=1.0\linewidth,
	boxrule=2.0pt, left=0.0pt, right=0.0pt, top=3pt, bottom=0.0pt,
	colbacktitle=docColor!20, coltitle=black, enhanced, attach boxed title to top left={xshift=3mm, yshift=-1mm},
	title=#2, #1}

\hypersetup{
	colorlinks=true,
	citecolor=blue,
	filecolor=magenta,
	linkcolor=blue,
	urlcolor=blue,
	pdftitle={SeminarAnnouncements},
	%bookmarks=true,,
}

\graphicspath{{./figs/}}

\definecolor{dgreen}{rgb}{0.0, 0.5, 0.0}

\allowdisplaybreaks

\newcommand*\mystrut[1]{\vrule width0pt height0pt depth#1\relax}
\renewcommand{\thesection}{\arabic{section}.}
\renewcommand{\thesubsection}{\arabic{section}.\arabic{subsection}}
\renewcommand{\thesubsubsection}{\arabic{section}.\arabic{subsection}.\arabic{subsubsection}}
\titleformat{\section}{\sffamily\large\bfseries}{\thesection}{1.0em}{}
\titleformat{\subsection}{\sffamily\large\bfseries}{\thesubsection}{1.0em}{}
\titleformat{\subsubsection}{\sffamily\normalsize\bfseries}{\thesubsubsection}{1.0em}{}
\titlespacing\section{0pt}{0pt plus 2pt minus 2pt}{0pt plus 2pt minus 2pt}
\renewcommand{\baselinestretch}{1.50}

\input{format/defs}

\fancypagestyle{titlepage}{
	\fancyhf{}
	\insertLXheader
}

\fancypagestyle{plain}{%
	\fancyhf{}
	\renewcommand{\headrulewidth}{0pt}
	\renewcommand{\footrulewidth}{0pt}
	\insertnormalpage
}

\pagestyle{plain}

\newcommand\TitleBackground{%
	\AddToShipoutPictureBG*{
		\begin{tikzpicture}[remember picture,overlay]
			\node[opacity=0.20, anchor=center, rotate=-90] 
			at (current page.center) [xshift=0em, yshift=0em]
			{\includegraphics[height=1.5\paperheight,keepaspectratio]{/Users/kannan/Library/CloudStorage/Dropbox/Kannan/0-templates/00-logos/bg-img/lms/lms-bg-1-portrait.pdf}};
		\end{tikzpicture}
	}
}

\newcommand{\lmsext}[2]{
	\textcolor{docColor}{#1} & \textcolor{docColor}{#2} \\
}

\newcommand{\lmsint}[2]{
	\rowcolor{cyan}
	\textcolor{docColor}{#1} & \textcolor{docColor}{[lms internal]~#2} \\
}

\newcommand{\meca}[2]{
	\rowcolor{docColor!99}
	\textcolor{white}{#1} & \textcolor{white}{[meca]~#2} \\
}

\setlength\arrayrulewidth{2pt}

\begin{document}
	\sffamily
	\thispagestyle{titlepage}
	\TitleBackground
	\vspace*{-3em}
	\begin{center}
		\huge \textbf{Seminars 2026 -- 27}
	\end{center}
	\vspace*{-3em}
	\begin{center}
		\Large 
		\begin{tabular}{| p{3.5cm} p{14cm} |}
			\hline

			\lmsext{September 03}{Nicolas Auffray (Institut Jean le Rond d'Alembert, Sorbonne)} 
			
			\meca{September 10}{Pole M{\'e}canique journ{\'e}e scientifique}
			
			\lmsint{September 17}{Fabrication capabilities for experiments (Vignesh)}
			
			\lmsext{September 24}{Alexander Erlich (University of Grenoble Alpes)}
			
			\lmsint{October 1}{}
			
			\lmsext{October 8}{Christelle Combescure (Saint-Cyr Co{\"e}tquidan Military Academy)}
			
			\lmsext{October 15}{Gowtham Nimmal Haribabu (TU Delft)}
			
			\meca{November 05}{}
			
			\lmsext{November 12}{Sathiskumar Ponnusami (City University of London)}
			
			\lmsext{November 19}{Jorge Grasa (University of Zaragoza)}
			
			\lmsext{November 26}{Adrian Buganza Tepole (Columbia University)}
			
			\meca{December 03}{}
			
			\lmsext{December 10}{}

			\lmsint{December 17}{}
			
			\hline
		\end{tabular}
	\end{center}
	
	\begin{center}
		\Large
		\begin{tabular}{p{4cm} p{15cm} }
			\raggedleft\textbf{Day and time:~} & Every Thursday, 1400 -- 1500 hrs (unless otherwise specified)\\ 
			\raggedleft\textbf{Venue:~} & Amphithéâtre Pole Mecanique\\
		\end{tabular}
		\vfill
		\textbullet \hspace*{0.5em} \href{https://calendar.google.com/calendar/embed?src=lms.seminaires%40gmail.com&ctz=Europe%2FZurich}{Google calendar link} \hspace*{0.5em} \textbullet \hspace*{0.5em} \href{https://polemecanique.github.io/lms-meca-seminars/}{Seminar webpage link} \hspace*{0.5em} \textbullet
	\end{center}
	
	
 \end{document}